% Options for packages loaded elsewhere
\PassOptionsToPackage{unicode}{hyperref}
\PassOptionsToPackage{hyphens}{url}
\PassOptionsToPackage{dvipsnames,svgnames,x11names}{xcolor}
%
\documentclass[
  10pt,
  a4paper,
]{article}
\usepackage{amsmath,amssymb}
\usepackage{iftex}
\ifPDFTeX
  \usepackage[T1]{fontenc}
  \usepackage[utf8]{inputenc}
  \usepackage{textcomp} % provide euro and other symbols
\else % if luatex or xetex
  \usepackage{unicode-math} % this also loads fontspec
  \defaultfontfeatures{Scale=MatchLowercase}
  \defaultfontfeatures[\rmfamily]{Ligatures=TeX,Scale=1}
\fi
\usepackage{lmodern}
\ifPDFTeX\else
  % xetex/luatex font selection
  \setmainfont[]{Noto Serif CJK JP}
  \setsansfont[]{Noto Sans CJK JP}
  \setmonofont[]{Noto Sans Mono CJK JP}
  \setmathfont[]{Latin Modern Math}
\fi
% Use upquote if available, for straight quotes in verbatim environments
\IfFileExists{upquote.sty}{\usepackage{upquote}}{}
\IfFileExists{microtype.sty}{% use microtype if available
  \usepackage[]{microtype}
  \UseMicrotypeSet[protrusion]{basicmath} % disable protrusion for tt fonts
}{}
\makeatletter
\@ifundefined{KOMAClassName}{% if non-KOMA class
  \IfFileExists{parskip.sty}{%
    \usepackage{parskip}
  }{% else
    \setlength{\parindent}{0pt}
    \setlength{\parskip}{6pt plus 2pt minus 1pt}}
}{% if KOMA class
  \KOMAoptions{parskip=half}}
\makeatother
\usepackage{xcolor}
\usepackage[margin=20mm]{geometry}
\usepackage{longtable,booktabs,array}
\usepackage{calc} % for calculating minipage widths
% Correct order of tables after \paragraph or \subparagraph
\usepackage{etoolbox}
\makeatletter
\patchcmd\longtable{\par}{\if@noskipsec\mbox{}\fi\par}{}{}
\makeatother
% Allow footnotes in longtable head/foot
\IfFileExists{footnotehyper.sty}{\usepackage{footnotehyper}}{\usepackage{footnote}}
\makesavenoteenv{longtable}
\setlength{\emergencystretch}{3em} % prevent overfull lines
\providecommand{\tightlist}{%
  \setlength{\itemsep}{0pt}\setlength{\parskip}{0pt}}
\setcounter{secnumdepth}{-\maxdimen} % remove section numbering
% XeLaTeX 用プリアンブル（pandoc の header-includes として読み込む）
% 日本語専用パッケージ（xeCJK / luatexja）なしで組めるよう、
% 本文フォントに和欧両対応の Noto Serif CJK JP を使い、
% 和文の行分割は XeTeX 組み込みのロケール改行に任せる。
\XeTeXlinebreaklocale "ja"
\XeTeXlinebreakskip=0pt plus 0.5pt
\XeTeXlinebreakpenalty=0
\usepackage{amsmath,amssymb}
\setlength{\parindent}{1em}
\setlength{\parskip}{0.4\baselineskip}
\renewcommand{\arraystretch}{1.25}
\usepackage{etoolbox}
\AtBeginEnvironment{longtable}{\small}

% Noto Serif CJK JP に無い記号（上付き・下付きの添字文字、ℝ・ℕ、⟨ ⟩ など）用の補助フォント。
% 該当箇所は latex/fallback-glyphs.lua が {\fallbackfont ...} で包む。
\newfontfamily\fallbackfont{DejaVu Serif}
\ifLuaTeX
  \usepackage{selnolig}  % disable illegal ligatures
\fi
\IfFileExists{bookmark.sty}{\usepackage{bookmark}}{\usepackage{hyperref}}
\IfFileExists{xurl.sty}{\usepackage{xurl}}{} % add URL line breaks if available
\urlstyle{same}
\hypersetup{
  colorlinks=true,
  linkcolor={blue},
  filecolor={Maroon},
  citecolor={Blue},
  urlcolor={blue},
  pdfcreator={LaTeX via pandoc}}

\author{}
\date{}

\begin{document}

\hypertarget{ux60c5ux5831ux5e7eux4f55ux5b66-ux4eacux5927ux7269ux7406ux5f62ux5f0fux30c6ux30adux30b9ux30c8}{%
\section{情報幾何学
京大物理形式テキスト}\label{ux60c5ux5831ux5e7eux4f55ux5b66-ux4eacux5927ux7269ux7406ux5f62ux5f0fux30c6ux30adux30b9ux30c8}}

作成日：2026-11-09

\hypertarget{ux4f7fux3044ux65b9}{%
\subsection{使い方}\label{ux4f7fux3044ux65b9}}

情報幾何学の基礎を8テーマに分け、京大物理と同じ「誘導つきの長文＋空欄補充＋選択＋問」形式で1テーマ＝大問1題にまとめた。確率分布の集まりを「曲がった空間」とみなし、計量（フィッシャー情報）、距離に似た量（ダイバージェンス）、2種類の直線（e-測地線と
m-測地線）を順に導入して、最尤推定や em
アルゴリズムを幾何の言葉で理解するところまで進む。各大問は30〜40分が目安。

\begin{itemize}
\tightlist
\item
  【ア】などの空欄には式・数値・語句を入れる。前の空欄の答えを後で使うことがある。
\item
  \{ ケ \} などの選択欄には①〜③から1つ選ぶ。
\item
  「問」は導出・証明・計算。論理を短く書く。
\item
  記号は京大物理にならい、カタカナ（Ⅰ・Ⅳ・Ⅶ）、イロハ（Ⅱ・Ⅴ・Ⅷ）、ひらがな（Ⅲ・Ⅵ）を順に使う。
\item
  解答と解説は末尾に2つに分けてまとめた。
\end{itemize}

前提知識：確率変数・期待値・分散（測度論的確率論テキストの問題Ⅲ・Ⅳ程度）、リーマン計量・クリストッフェル記号・ガウス曲率（微分幾何学テキストの問題Ⅱ〜Ⅳ程度）、凸関数とルジャンドル変換の初歩。

\begin{longtable}[]{@{}
  >{\raggedright\arraybackslash}p{(\columnwidth - 4\tabcolsep) * \real{0.3333}}
  >{\raggedright\arraybackslash}p{(\columnwidth - 4\tabcolsep) * \real{0.3333}}
  >{\raggedright\arraybackslash}p{(\columnwidth - 4\tabcolsep) * \real{0.3333}}@{}}
\toprule\noalign{}
\begin{minipage}[b]{\linewidth}\raggedright
大問
\end{minipage} & \begin{minipage}[b]{\linewidth}\raggedright
テーマ
\end{minipage} & \begin{minipage}[b]{\linewidth}\raggedright
到達目標
\end{minipage} \\
\midrule\noalign{}
\endhead
\bottomrule\noalign{}
\endlastfoot
Ⅰ & 統計モデルとフィッシャー情報計量 &
スコア関数とフィッシャー情報を計算し、計量として扱える \\
Ⅱ & KL・f-ダイバージェンス &
ダイバージェンスの局所展開がフィッシャー計量になることを示せる \\
Ⅲ & 指数型分布族と双対座標 & θ 座標・η 座標・ポテンシャル ψ, φ
とルジャンドル変換 \\
Ⅳ & α-接続と双対接続 & e-接続・m-接続、双対性、平坦性の判定 \\
Ⅴ & 双対平坦空間と拡張ピタゴラスの定理 &
ブレグマン・ダイバージェンスと直交分解 \\
Ⅵ & 射影定理と統計的推論 & 最尤推定＝m-射影、最大エントロピー、em
アルゴリズム \\
Ⅶ & クラメール＝ラオの不等式 &
推定精度の限界とフィッシャー情報の関係、有効推定量の条件 \\
Ⅷ & 正規分布族の幾何 & 双曲平面としての構造、曲率・測地距離の計算 \\
\end{longtable}

共通の記号：統計モデル \{p(x; ξ)\} の対数尤度を l(x; ξ) = log p(x;
ξ)、∂{\fallbackfont ᵢ} = ∂/∂ξ{\fallbackfont ⁱ} とし、期待値 E
は特に断らない限り p(x; ξ)
に関するものとする。微分と積分の順序交換は自由に行ってよい。

\hypertarget{ux554fux984cux2170-ux7d71ux8a08ux30e2ux30c7ux30ebux3068ux30d5ux30a3ux30c3ux30b7ux30e3ux30fcux60c5ux5831ux8a08ux91cf}{%
\subsection{問題Ⅰ　統計モデルとフィッシャー情報計量}\label{ux554fux984cux2170-ux7d71ux8a08ux30e2ux30c7ux30ebux3068ux30d5ux30a3ux30c3ux30b7ux30e3ux30fcux60c5ux5831ux8a08ux91cf}}

次の文章を読んで、【　】に適した式・数値を、\{　\}
からは適切なものを一つ選びその番号を、それぞれ答えよ。また、問1・問2では指示にしたがって答えよ。

\begin{enumerate}
\def\labelenumi{(\arabic{enumi})}
\tightlist
\item
  n 個の実パラメータ ξ = (ξ¹, \ldots, ξ{\fallbackfont ⁿ})
  で滑らかに指定される確率分布の族 S = \{p(x; ξ)\} を統計モデルという。ξ
  を座標とみなせば、S は n 次元の多様体になる。各点 ξ
  で、対数尤度の偏微分 ∂{\fallbackfont ᵢ}l をスコア関数という。∫ p(x; ξ)
  dx = 1 の両辺を ξ{\fallbackfont ⁱ} で微分すると
\end{enumerate}

\[E[\partial_i l] = \int \partial_i p(x;\xi)\, dx = \text{【ア】}\]

となり、スコアの平均は 0 である。そこで、スコアの共分散

\[g_{ij}(\xi) = E[\partial_i l\; \partial_j l]\]

をフィッシャー情報行列という。任意の実数 c{\fallbackfont ₁}, \ldots,
c{\fallbackfont ₙ} に対し Σ
c{\fallbackfont ᵢ}c{\fallbackfont ⱼ}g{\fallbackfont ᵢⱼ} = E{[}(Σ
c{\fallbackfont ᵢ}∂{\fallbackfont ᵢ}l)²{]} ≥ 0 なので g
は半正定値であり、スコアが一次独立なら正定値になる。こうして g は S
上のリーマン計量（フィッシャー計量）を定める。

また、∂{\fallbackfont ᵢ}∂{\fallbackfont ⱼ}l =
(∂{\fallbackfont ᵢ}∂{\fallbackfont ⱼ}p)/p − ∂{\fallbackfont ᵢ}l
∂{\fallbackfont ⱼ}l と E{[}(∂{\fallbackfont ᵢ}∂{\fallbackfont ⱼ}p)/p{]}
= 【イ】 から、g{\fallbackfont ᵢⱼ} =
−E{[}∂{\fallbackfont ᵢ}∂{\fallbackfont ⱼ}l{]}
とも書ける（問1）。対数尤度の「山のとがり具合」の平均がフィッシャー情報である。

\begin{enumerate}
\def\labelenumi{(\arabic{enumi})}
\setcounter{enumi}{1}
\tightlist
\item
  具体例を計算しよう。
\end{enumerate}

\begin{itemize}
\tightlist
\item
  ベルヌーイ分布 p(x; θ) = θ{\fallbackfont ˣ}(1 −
  θ)¹{\fallbackfont ⁻ˣ}（x = 0, 1、0 \textless{} θ \textless{} 1）では
  ∂l = (x − θ)/(θ(1 − θ)) なので g(θ) = 【ウ】
\item
  分散 σ² が既知の正規分布 N(μ, σ²) では g(μ) = 【エ】
\item
  平均・標準偏差の両方を座標とする正規分布族 (μ, σ) では g = diag(1/σ²,
  【オ】)
\end{itemize}

ベルヌーイの例では θ が 0 や 1 に近いほど g が大きい。θ = 0.01 と 0.02
は、θ = 0.50 と 0.51 よりも、データから見分けやすいからである。

\begin{enumerate}
\def\labelenumi{(\arabic{enumi})}
\setcounter{enumi}{2}
\tightlist
\item
  座標を ξ から ζ に取り替えると、フィッシャー情報は
\end{enumerate}

\[\tilde g_{ab}(\zeta) = \sum_{i,j} \frac{\partial \xi^i}{\partial \zeta^a} \frac{\partial \xi^j}{\partial \zeta^b}\, g_{ij}(\xi)\]

と、リーマン計量と同じ規則で変換される。ベルヌーイ分布で θ =
sin²{\fallbackfont ϕ}（0 \textless{} {\fallbackfont ϕ} \textless{}
π/2）とおくと {\fallbackfont g̃}({\fallbackfont ϕ}) = 【カ】
と定数になる。したがってベルヌーイ族全体は長さ 【キ】
の線分と等長である。

より一般に、k + 1
個の値をとる離散分布全体（確率単体）のフィッシャー計量は ds² = Σ
(dp{\fallbackfont ᵢ})²/p{\fallbackfont ᵢ} = 4 Σ (d√p{\fallbackfont ᵢ})²
と書けるので、y{\fallbackfont ᵢ} = √p{\fallbackfont ᵢ} と変換すれば Σ
y{\fallbackfont ᵢ}² = 1
を満たす単位球面の一部に4倍の計量を入れたもの、すなわち半径 【ク】
の球面の一部と等長になる。

フィッシャー計量は、十分統計量による変換（情報を失わないデータの要約）で不変である。この性質について、チェンツォフの定理によれば、十分統計量による変換で不変な確率単体上のリーマン計量は
\{ ケ \}。

\begin{itemize}
\tightlist
\item
  ① 定数倍を除いてフィッシャー計量に限る
\item
  ② 無数に存在する
\item
  ③ 存在しない
\end{itemize}

問1　g{\fallbackfont ᵢⱼ} = −E{[}∂{\fallbackfont ᵢ}∂{\fallbackfont ⱼ}l{]}
を導け。

問2　ポアソン分布 p(x; λ) = e\^{}(−λ) λ{\fallbackfont ˣ}/x!（x = 0, 1,
2, \ldots）のフィッシャー情報 g(λ) を求めよ。また、λ = ν²
と座標変換すると計量が定数になることを示せ。

\hypertarget{ux554fux984cux2171-klux30c0ux30a4ux30d0ux30fcux30b8ux30a7ux30f3ux30b9ux3068-f-ux30c0ux30a4ux30d0ux30fcux30b8ux30a7ux30f3ux30b9}{%
\subsection{問題Ⅱ　KLダイバージェンスと
f-ダイバージェンス}\label{ux554fux984cux2171-klux30c0ux30a4ux30d0ux30fcux30b8ux30a7ux30f3ux30b9ux3068-f-ux30c0ux30a4ux30d0ux30fcux30b8ux30a7ux30f3ux30b9}}

次の文章を読んで、【　】に適した式・数値を、\{　\}
からは適切なものを一つ選びその番号を、それぞれ答えよ。また、問1・問2では指示にしたがって答えよ。

\begin{enumerate}
\def\labelenumi{(\arabic{enumi})}
\tightlist
\item
  2つの確率分布 p, q
  の「隔たり」を測る量として、カルバック＝ライブラー（KL）ダイバージェンス
\end{enumerate}

\[D(p \,\|\, q) = \int p(x) \log \frac{p(x)}{q(x)}\, dx\]

が基本的である。不等式 log t ≤ t − 1（等号は t = 1）を t = q/p
に使うと、−D(p‖q) = ∫ p log(q/p) dx ≤ ∫ p (q/p − 1) dx = 【イ】
となり、D(p‖q) ≥ 0 が従う（ギブスの不等式）。等号は p = q
のときに限る。ただし一般に D(p‖q) ≠ D(q‖p) であり、KL ダイバージェンスは
\{ ロ \}。

\begin{itemize}
\tightlist
\item
  ① 非対称で、三角不等式も一般には満たさないので、距離ではない
\item
  ② 対称だが、三角不等式を満たさない
\item
  ③ 距離の公理をすべて満たす
\end{itemize}

\begin{enumerate}
\def\labelenumi{(\arabic{enumi})}
\setcounter{enumi}{1}
\tightlist
\item
  距離ではないが、近い2点の間では KL
  ダイバージェンスはフィッシャー計量と結びつく。統計モデル上の2点 ξ, ξ′
  について f(ξ′) = D(p(·; ξ) ‖ p(·; ξ′)) とおくと f(ξ) = 0 である。ξ′
  で微分すると ∂f/∂ξ′{\fallbackfont ⁱ} = −E{[}∂{\fallbackfont ᵢ}l(x;
  ξ′){]}（期待値は p(x; ξ) に関する）なので、ξ′ = ξ では 【ハ】
  となる。さらにもう1回微分して ξ′ = ξ とおくと
  −E{[}∂{\fallbackfont ᵢ}∂{\fallbackfont ⱼ}l{]} = 【ニ】 となるから
\end{enumerate}

\[D\big(p_\xi \,\|\, p_{\xi + d\xi}\big) = \frac{1}{2} \sum_{i,j} g_{ij}(\xi)\, d\xi^i d\xi^j + O(|d\xi|^3)\]

を得る。KL
ダイバージェンスは、局所的にはフィッシャー計量で測った距離の2乗の半分である。

例として、分散が共通の正規分布では D(N(μ{\fallbackfont ₁}, σ²) ‖
N(μ{\fallbackfont ₂}, σ²)) = 【ホ】 と正確に計算でき、問題Ⅰの g(μ) =
1/σ² と整合する。ポアソン分布では D(Po(λ{\fallbackfont ₁}) ‖
Po(λ{\fallbackfont ₂})) = 【ヘ】 であり、λ{\fallbackfont ₂} =
λ{\fallbackfont ₁} + dλ と展開すると D ≈ (1/2)(【ト】)dλ² となる。

\begin{enumerate}
\def\labelenumi{(\arabic{enumi})}
\setcounter{enumi}{2}
\tightlist
\item
  凸関数 f（f(1) = 0）を用いて
\end{enumerate}

\[D_f(p \,\|\, q) = \int p(x)\, f\!\left(\frac{q(x)}{p(x)}\right) dx\]

と定めたものを f-ダイバージェンスという。イェンセンの不等式から D\_f ≥ 0
である。いくつかの例をあげる。

\begin{longtable}[]{@{}lll@{}}
\toprule\noalign{}
名前 & f(t) & f″(1) \\
\midrule\noalign{}
\endhead
\bottomrule\noalign{}
\endlastfoot
KL ダイバージェンス D(p‖q) & 【チ】 & 【リ】 \\
逆向き KL D(q‖p) & t log t & 1 \\
ヘリンジャー距離の2乗 ∫(√p − √q)² dx & (√t − 1)² & 【ヌ】 \\
全変動距離 & \textbar t − 1\textbar/2 & 定義されない（t = 1
で微分不可能） \\
\end{longtable}

f が t = 1 で2回微分可能なら、(2) と同様の計算で、近い2点 ξ と ξ + dξ
の間の f-ダイバージェンスは (f″(1)/2) Σ g{\fallbackfont ᵢⱼ}
dξ{\fallbackfont ⁱ} dξ{\fallbackfont ʲ} で近似される。どの滑らかな
f-ダイバージェンスも、局所的にはフィッシャー計量の定数倍に帰着するのである。

f-ダイバージェンスは、データを統計量 T
で変換すると増えない（データ処理不等式
D\_f(p{\fallbackfont ∘}T{\fallbackfont ⁻}¹ ‖
q{\fallbackfont ∘}T{\fallbackfont ⁻}¹) ≤ D\_f(p‖q)）。T
が十分統計量なら等号が成り立つ。問題Ⅰで述べたフィッシャー計量の不変性は、この性質の局所版である。

問1　ギブスの不等式 D(p‖q) ≥ 0 を、等号成立条件も含めて示せ。

問2　D(N(μ{\fallbackfont ₁}, σ²) ‖ N(μ{\fallbackfont ₂}, σ²))
を直接計算して【ホ】を確かめよ。

\hypertarget{ux554fux984cux2172-ux6307ux6570ux578bux5206ux5e03ux65cfux3068ux53ccux5bfeux5ea7ux6a19}{%
\subsection{問題Ⅲ　指数型分布族と双対座標}\label{ux554fux984cux2172-ux6307ux6570ux578bux5206ux5e03ux65cfux3068ux53ccux5bfeux5ea7ux6a19}}

次の文章を読んで、【　】に適した式・数値・語句を、\{　\}
からは適切なものを一つ選びその番号を、それぞれ答えよ。また、問1・問2では指示にしたがって答えよ。

\begin{enumerate}
\def\labelenumi{(\arabic{enumi})}
\tightlist
\item
  統計量 F{\fallbackfont ₁}(x), \ldots, F{\fallbackfont ₙ}(x)
  とパラメータ θ = (θ¹, \ldots, θ{\fallbackfont ⁿ}) を用いて
\end{enumerate}

\[p(x;\theta) = \exp\Big( \sum_{i=1}^{n} \theta^i F_i(x) - \psi(\theta) \Big), \qquad \psi(\theta) = \log \int \exp\Big( \sum_i \theta^i F_i(x) \Big) dx\]

と書ける分布の族を指数型分布族、θ を自然パラメータ（θ
座標）という（積分は基準となる測度に関するものとする）。ψ は確率の総和を
1 にする正規化項で、対数分配関数あるいはキュムラント母関数と呼ばれる。

ψ を微分すると ∂{\fallbackfont ᵢ}ψ =
E{[}F{\fallbackfont ᵢ}{]}、∂{\fallbackfont ᵢ}∂{\fallbackfont ⱼ}ψ =
Cov(F{\fallbackfont ᵢ}, F{\fallbackfont ⱼ})
となる（問1）。一方、対数尤度 l = Σ θ{\fallbackfont ⁱ}F{\fallbackfont ᵢ}
− ψ(θ) の2階微分は ∂{\fallbackfont ᵢ}∂{\fallbackfont ⱼ}l =
−∂{\fallbackfont ᵢ}∂{\fallbackfont ⱼ}ψ で x によらないので、問題Ⅰより
Cov(F{\fallbackfont ᵢ}, F{\fallbackfont ⱼ}) = 【あ】
が成り立つ。共分散行列は半正定値なので、ψ は凸関数である。

\begin{enumerate}
\def\labelenumi{(\arabic{enumi})}
\setcounter{enumi}{1}
\tightlist
\item
  ベルヌーイ分布は p(x) = exp(θx − ψ(θ))、θ = log(p/(1 −
  p))（ロジット）と書けて、ψ(θ) = 【い】 である。期待値は η = ψ′(θ) =
  【う】 となり、これは成功確率 p そのものである。θ
  座標でのフィッシャー情報は g(θ) = ψ″(θ) = 【え】 である。問題Ⅰで求めた
  p 座標での値 1/(p(1 − p)) と見かけは異なるが、(dp/dθ)² = (p(1 − p))²
  を掛ければ一致し、座標変換則と整合している。
\end{enumerate}

正規分布 N(μ, σ²) は F{\fallbackfont ₁} = x、F{\fallbackfont ₂} = x²
として指数型分布族になり、θ¹ = μ/σ²、θ² = 【お】
である。対応する期待値は η{\fallbackfont ₁} = E{[}x{]} =
μ、η{\fallbackfont ₂} = E{[}x²{]} = 【か】 である。

\begin{enumerate}
\def\labelenumi{(\arabic{enumi})}
\setcounter{enumi}{2}
\tightlist
\item
  η{\fallbackfont ᵢ} = ∂{\fallbackfont ᵢ}ψ(θ) で定まる η =
  (η{\fallbackfont ₁}, \ldots, η{\fallbackfont ₙ}) を期待値パラメータ（η
  座標）という。ψ が狭義凸なら θ {\fallbackfont ↦} η は1対1なので、η
  も座標になる。θ と η を結ぶのがルジャンドル変換である。
\end{enumerate}

\[\varphi(\eta) = \sup_{\theta} \big\{ \theta \cdot \eta - \psi(\theta) \big\}\]

上限は ∇ψ(θ) = η を満たす θ で達成され、φ(η) = θ · η − ψ(θ) = E{[}l{]}
となる。つまり φ(η) は分布の 【き】 に等しい。対応する点では ψ(θ) + φ(η)
− θ · η = 【く】 が成り立ち、さらに θ{\fallbackfont ⁱ} =
∂φ/∂η{\fallbackfont ᵢ} と、逆向きの関係も同じ形で成り立つ。ヤコビ行列
∂η{\fallbackfont ᵢ}/∂θ{\fallbackfont ʲ} = g{\fallbackfont ᵢⱼ} と
∂θ{\fallbackfont ⁱ}/∂η{\fallbackfont ⱼ} は互いに逆行列であり、η
座標でのフィッシャー計量は ψ のヘッセ行列の逆行列、すなわち φ
のヘッセ行列になる。ベルヌーイ分布では φ(η) = 【け】 である。

こうして指数型分布族には、凸関数 ψ と φ に支えられた2つの座標 θ, η
が備わる。この双対構造が以降の大問の主役である。最初の応用として、データ
x{\fallbackfont ₁}, \ldots, x\_N から θ を最尤推定すると、推定値
{\fallbackfont θ̂} は \{ こ \}。

\begin{itemize}
\tightlist
\item
  ① 統計量の標本平均と期待値パラメータが一致する条件
  η({\fallbackfont θ̂}) = (1/N) Σ{\fallbackfont ₖ} F(x{\fallbackfont ₖ})
  で決まる
\item
  ② 尤度方程式に解をもたないので数値的にしか求まらない
\item
  ③ 統計量の標本中央値で決まる
\end{itemize}

問1　∂{\fallbackfont ᵢ}ψ = E{[}F{\fallbackfont ᵢ}{]} と
∂{\fallbackfont ᵢ}∂{\fallbackfont ⱼ}ψ = Cov(F{\fallbackfont ᵢ},
F{\fallbackfont ⱼ}) を示せ。

問2　ポアソン分布を指数型分布族として表し、θ, ψ(θ), η, θ
座標でのフィッシャー情報、φ(η) を求めよ。

\hypertarget{ux554fux984cux2173-ux3b1-ux63a5ux7d9aux3068ux53ccux5bfeux63a5ux7d9a}{%
\subsection{問題Ⅳ　α-接続と双対接続}\label{ux554fux984cux2173-ux3b1-ux63a5ux7d9aux3068ux53ccux5bfeux63a5ux7d9a}}

次の文章を読んで、【　】に適した式・数値・語句を、\{　\}
からは適切なものを一つ選びその番号を、それぞれ答えよ。また、問1・問2では指示にしたがって答えよ。

\begin{enumerate}
\def\labelenumi{(\arabic{enumi})}
\tightlist
\item
  多様体上で「まっすぐ」を決めるのがアファイン接続 ∇
  である。座標基底に対する係数 Γ{\fallbackfont ᵢⱼ},{\fallbackfont ₖ} =
  {\fallbackfont ⟨}∇(∂{\fallbackfont ᵢ}) ∂{\fallbackfont ⱼ},
  ∂{\fallbackfont ₖ⟩}（{\fallbackfont ⟨} , {\fallbackfont ⟩}
  はフィッシャー計量）で接続を表す。微分幾何学で学んだレヴィ＝チヴィタ接続はその一つにすぎず、統計モデルには実数
  α でパラメータ付けされた接続の族（α-接続）が自然に備わっている。
\end{enumerate}

\[\Gamma^{(\alpha)}_{ij,k} = E\Big[ \Big( \partial_i \partial_j l + \frac{1 - \alpha}{2}\, \partial_i l\, \partial_j l \Big) \partial_k l \Big], \qquad T_{ijk} = E[\partial_i l\, \partial_j l\, \partial_k l]\]

α = 1 を e-接続（指数接続）、α = −1 を m-接続（混合接続）、α = 0
をレヴィ＝チヴィタ接続という。3階対称テンソル T を用いると
Γ\^{}(α){\fallbackfont ᵢⱼ},{\fallbackfont ₖ} =
Γ\^{}(0){\fallbackfont ᵢⱼ},{\fallbackfont ₖ} − 【ア】
T{\fallbackfont ᵢⱼₖ} と書け、α-接続はレヴィ＝チヴィタ接続を T
の方向にずらしたものである。

α-接続の間には次の関係がある。

\[\partial_k g_{ij} = \Gamma^{(\alpha)}_{ki,j} + \Gamma^{(\text{【イ】})}_{kj,i}\]

これは「ベクトル X を ∇ で、Y を ∇* で平行移動すると内積
{\fallbackfont ⟨}X, Y{\fallbackfont ⟩}
が保たれる」ことを意味し、この関係にある ∇ と ∇*
を互いに双対な接続という。とくに e-接続と
m-接続は双対である。自分自身と双対になるのは α = 【ウ】
のときだけで、これは計量を保つレヴィ＝チヴィタ接続の性質そのものである。

\begin{enumerate}
\def\labelenumi{(\arabic{enumi})}
\setcounter{enumi}{1}
\tightlist
\item
  指数型分布族を θ 座標で表すと ∂{\fallbackfont ᵢ}∂{\fallbackfont ⱼ}l =
  −∂{\fallbackfont ᵢ}∂{\fallbackfont ⱼ}ψ は x
  によらないので、E{[}∂{\fallbackfont ₖ}l{]} = 0 から
  Γ\^{}(e){\fallbackfont ᵢⱼ},{\fallbackfont ₖ} = 【エ】 となる。すなわち
  θ 座標が e-接続の直線座標（アファイン座標）であり、指数型分布族は
  e-平坦である。一方、分布 p{\fallbackfont ₀}, p{\fallbackfont ₁},
  \ldots, p{\fallbackfont ₙ} の混合
\end{enumerate}

\[p(x;\eta) = \sum_{i=1}^{n} \eta_i\, p_i(x) + \Big(1 - \sum_{i=1}^{n} \eta_i\Big) p_0(x)\]

からなる混合型分布族では ∂{\fallbackfont ᵢ}∂{\fallbackfont ⱼ}p = 0
であり、Γ\^{}(m){\fallbackfont ᵢⱼ},{\fallbackfont ₖ} = ∫
∂{\fallbackfont ᵢ}∂{\fallbackfont ⱼ}p ∂{\fallbackfont ₖ}l dx = 【オ】
となるので m-平坦である。双対な接続の一方の曲率が 0 なら、他方の曲率も 0
になることが知られており、指数型分布族は e-接続に関しても
m-接続に関しても 【カ】 である。m-接続の直線座標は η 座標である。

\begin{enumerate}
\def\labelenumi{(\arabic{enumi})}
\setcounter{enumi}{2}
\tightlist
\item
  2つの分布 p, q を結ぶ測地線を書き下そう。
\end{enumerate}

\begin{itemize}
\tightlist
\item
  m-測地線：p(t) = 【キ】（0 ≤ t ≤ 1）。確率を足し合わせる「混合」の直線
\item
  e-測地線：log p(t) = (1 − t) log p + t log q −
  c(t)。対数をとってから足し合わせる「幾何平均」の直線で、正規化定数は
  c(t) = 【ク】
\end{itemize}

同じ2点を結んでいても、2本の測地線は一般に異なる曲線である。

\begin{enumerate}
\def\labelenumi{(\arabic{enumi})}
\setcounter{enumi}{3}
\tightlist
\item
  ベルヌーイ分布の θ 座標では、T は ψ の3階微分に等しく T =
  ψ{\fallbackfont ‴}(θ) = 【ケ】 である。p = 1/2 では T = 0
  となり、この点では α によらず接続が一致する。
\end{enumerate}

なお、指数型分布族は
e-接続・m-接続に関して平坦だが、レヴィ＝チヴィタ接続に関しては \{ コ
\}。

\begin{itemize}
\tightlist
\item
  ① 一般には平坦でない（問題Ⅷの正規分布族が例）
\item
  ② 常に平坦である
\item
  ③ 接続そのものが定義できない
\end{itemize}

問1　e-接続と m-接続について、∂{\fallbackfont ₖ}g{\fallbackfont ᵢⱼ} =
Γ\^{}(e){\fallbackfont ₖᵢ},{\fallbackfont ⱼ} +
Γ\^{}(m){\fallbackfont ₖⱼ},{\fallbackfont ᵢ} を直接計算で確かめよ。

問2　N(μ{\fallbackfont ₁}, 1) と N(μ{\fallbackfont ₂}, 1) を結ぶ
e-測地線上の分布が、すべて分散 1 の正規分布 N((1 − t)μ{\fallbackfont ₁}
+ tμ{\fallbackfont ₂}, 1) になることを示せ。

\hypertarget{ux554fux984cux2174-ux53ccux5bfeux5e73ux5766ux7a7aux9593ux3068ux62e1ux5f35ux30d4ux30bfux30b4ux30e9ux30b9ux306eux5b9aux7406}{%
\subsection{問題Ⅴ　双対平坦空間と拡張ピタゴラスの定理}\label{ux554fux984cux2174-ux53ccux5bfeux5e73ux5766ux7a7aux9593ux3068ux62e1ux5f35ux30d4ux30bfux30b4ux30e9ux30b9ux306eux5b9aux7406}}

次の文章を読んで、【　】に適した式・語句を、\{　\}
からは適切なものを一つ選びその番号を、それぞれ答えよ。また、問1・問2では指示にしたがって答えよ。

\begin{enumerate}
\def\labelenumi{(\arabic{enumi})}
\tightlist
\item
  問題Ⅲ・Ⅳで見た構造を抽象化する。多様体 M 上に、互いに双対な接続 ∇, ∇*
  がともに平坦で、∇ の直線座標 θ と ∇* の直線座標 η が、凸関数 ψ(θ),
  φ(η) を用いて η = ∇ψ(θ)、θ = ∇φ(η)、ψ(θ) + φ(η) − θ · η = 0
  で結ばれているとき、M
  を双対平坦空間という。指数型分布族はその典型例である。
\end{enumerate}

2点 P, Q（座標をそれぞれ θ\_P, η\_P
などと書く）に対し、正準ダイバージェンスを

\[D[P:Q] = \psi(\theta_P) + \varphi(\eta_Q) - \theta_P \cdot \eta_Q\]

で定める。指数型分布族でこれを計算すると、分布 q（点 Q）から分布 p（点
P）への KL ダイバージェンスと一致し、D{[}P:Q{]} = 【イ】
となる（引数の順序が逆になることに注意）。

φ(η\_Q) = θ\_Q · η\_Q − ψ(θ\_Q) と η\_Q = ∇ψ(θ\_Q) を代入すると

\[D[P:Q] = \psi(\theta_P) - \psi(\theta_Q) - \text{【ロ】}\]

と書き直せる。これは「関数 ψ と、点 θ\_Q
での接平面（1次近似）との高さの差」であり、ブレグマン・ダイバージェンスと呼ばれる。ψ
の 【ハ】 から D{[}P:Q{]} ≥ 0 であり、等号は P = Q
のときに限る（問2）。たとえば ψ(θ) = \textbar θ\textbar²/2（分散 1
の正規分布の平均パラメータ）なら η = θ で、D{[}P:Q{]} = 【ニ】
となり、通常のユークリッド距離の2乗の半分に一致する。

\begin{enumerate}
\def\labelenumi{(\arabic{enumi})}
\setcounter{enumi}{1}
\tightlist
\item
  3点 P, Q, R に対し、定義に代入して整理すると
\end{enumerate}

\[D[P:R] = D[P:Q] + D[Q:R] - \text{【ホ】}\]

という「余弦定理」が成り立つ（問1）。最後の項は、Q から P へ向かう
∇-測地線（θ 座標で直線）の方向 θ\_P − θ\_Q と、Q から R へ向かう
∇*-測地線（η 座標で直線）の方向 η\_R − η\_Q
の内積である。一般に、接ベクトル X を θ 座標の成分 X{\fallbackfont ⁱ}
で、Y を η 座標の成分 Y{\fallbackfont ᵢ}
で表すと、フィッシャー計量による内積は {\fallbackfont ⟨}X,
Y{\fallbackfont ⟩} = 【ヘ】
と、計量を介さず単純な和になる。双対座標の最大の利点である。

したがって、Q において ∇-測地線 QP と ∇*-測地線 QR
が直交すれば、D{[}P:R{]} = D{[}P:Q{]} + D{[}Q:R{]}
が成り立つ。これを拡張ピタゴラスの定理という。KL
ダイバージェンスの言葉に直すと、e-測地線 QP と m-測地線 QR が Q
で直交するとき

\[D(r \,\|\, p) = D(q \,\|\, p) + \text{【ト】}\]

が成り立つ（p, q, r はそれぞれ点 P, Q, R
の分布）。この定理は、次の問題Ⅵで「最も近い点への射影」を特徴づける基礎になる。

なお、ダイバージェンスの引数を入れ替えた D{[}Q:P{]} は、一般に \{ チ
\}。

\begin{itemize}
\tightlist
\item
  ① D{[}P:Q{]} と異なり、入れ替えると ∇ と ∇*（θ と η、ψ と
  φ）の役割が入れ替わる
\item
  ② 常に D{[}P:Q{]} に等しい
\item
  ③ 負の値をとりうる
\end{itemize}

問1　正準ダイバージェンスの定義から、(2) の余弦定理を導け。

問2　ψ が狭義凸のとき、ブレグマン・ダイバージェンスが非負で、0
になるのは θ\_P = θ\_Q のときに限ることを示せ。

\hypertarget{ux554fux984cux2175-ux5c04ux5f71ux5b9aux7406ux6700ux5c24ux63a8ux5b9aux6700ux5927ux30a8ux30f3ux30c8ux30edux30d4ux30fcem-ux30a2ux30ebux30b4ux30eaux30baux30e0}{%
\subsection{問題Ⅵ　射影定理・最尤推定・最大エントロピー・em
アルゴリズム}\label{ux554fux984cux2175-ux5c04ux5f71ux5b9aux7406ux6700ux5c24ux63a8ux5b9aux6700ux5927ux30a8ux30f3ux30c8ux30edux30d4ux30fcem-ux30a2ux30ebux30b4ux30eaux30baux30e0}}

次の文章を読んで、【　】に適した式・語句を、\{　\}
からは適切なものを一つ選びその番号を、それぞれ答えよ。また、問1・問2では指示にしたがって答えよ。以下、標本空間は有限集合とする。

\begin{enumerate}
\def\labelenumi{(\arabic{enumi})}
\tightlist
\item
  データ x{\fallbackfont ₁}, \ldots, x\_N
  が得られたとき、各値の出現割合を確率とする経験分布 {\fallbackfont p̂}
  を考える。モデル M = \{p(x; θ)\} の分布との KL ダイバージェンスは
\end{enumerate}

\[D(\hat p \,\|\, p_\theta) = \sum_x \hat p(x) \log \hat p(x) - \text{【あ】}\]

と書ける。第1項は θ によらないので、D({\fallbackfont p̂}‖p\_θ)
を最小にする θ は最尤推定量に一致する。最尤推定は、データ点
{\fallbackfont p̂} からモデル M への「最も近い点」探しである。

M 上の点 q のうち D(p‖q) を最小にするものを求める操作を、p から M への
【い】 という。問題Ⅴの拡張ピタゴラスの定理より、{\fallbackfont q̂}
が最小点ならば、p から {\fallbackfont q̂} への m-測地線は
{\fallbackfont q̂} で M と直交する（射影定理）。さらに M が 【う】
であれば、そのような点はただ1つに定まる。逆向きの D(q‖p) を q
について最小にする操作は e-射影と呼ばれ、射影先の集合が
m-平坦なら一意である。

M が指数型分布族なら、問題Ⅲで見たとおり最尤推定量は η({\fallbackfont θ̂})
= (1/N) Σ{\fallbackfont ₖ} F(x{\fallbackfont ₖ})
を満たす。たとえばベルヌーイ分布で N 回中 k
回成功したとき、成功確率の推定値は 【え】 である。幾何的には、「F
の期待値が標本平均に等しい分布」全体の集合 A は m-平坦で
{\fallbackfont p̂} を含み、e-平坦な M と点 {\fallbackfont q̂}
でちょうど直交して交わる。

\begin{enumerate}
\def\labelenumi{(\arabic{enumi})}
\setcounter{enumi}{1}
\tightlist
\item
  最大エントロピー原理。期待値 E{[}F{]} = c を満たす分布全体の集合 A
  の中で、エントロピー H(p) = −Σ p log p
  が最大のものを求めよう。標本空間の要素数を K、一様分布を u とすると
  D(p‖u) = log K − H(p) なので、エントロピー最大化は u からの D(p‖u)
  の最小化、すなわち u から A への
  e-射影である。ラグランジュの未定乗数法で解くと、最大点は p*(x) =
  【お】 の形の指数型分布族の一員になる。ここで θ は制約 E{[}F{]} = c
  に対応する 【か】 である。
\end{enumerate}

拡張ピタゴラスの定理を使うと、A の任意の分布 p について D(p‖u) = D(p‖p*)
+ D(p*‖u) が成り立ち、H(p) ≤ H(p*)
が直ちに従う。統計力学のカノニカル分布 e\^{}(−βE)/Z は、F
をエネルギー、θ = −β とした特別な場合である。

\begin{enumerate}
\def\labelenumi{(\arabic{enumi})}
\setcounter{enumi}{2}
\tightlist
\item
  em アルゴリズム。観測できる変数 x と隠れ変数 z をもつモデル M = \{p(x,
  z; θ)\}（e-平坦とする）を考える。(x, z) の同時分布のうち、x
  の周辺分布が経験分布 {\fallbackfont p̂}(x)
  に一致するもの全体をデータ多様体 D とよぶと、D は m-平坦である。点 p ∈
  D と点 q ∈ M の間のダイバージェンス D(p‖q) を交互に最小化する。
\end{enumerate}

\begin{itemize}
\tightlist
\item
  e ステップ：q を固定して p ∈ D について最小化する。これは q から D
  への 【き】 であり、解は p(x, z) = {\fallbackfont p̂}(x) q(z \textbar{}
  x)（隠れ変数の条件付き分布で補完）となる
\item
  m ステップ：p を固定して q ∈ M について最小化する。これは p から M
  への m-射影、すなわち補完した完全データに対する最尤推定である
\end{itemize}

各ステップで D(p‖q) は 【く】
なので、アルゴリズムは必ず収束する方向に進むが、到達するのは一般に
【け】 であり、大域的な最小点とは限らない。統計学で広く使われる EM
アルゴリズムとの関係については、\{ こ \}。

\begin{itemize}
\tightlist
\item
  ① M が指数型分布族なら、em アルゴリズムと EM アルゴリズムは一致する
\item
  ② 2つは常に異なる推定値を与える
\item
  ③ EM アルゴリズムは乱数を使う確率的手法なので比較できない
\end{itemize}

問1　(1) の KL ダイバージェンスの式を導き、最尤推定が {\fallbackfont p̂}
から M への射影であることを説明せよ。

問2　サイコロの目 \{1, 2, \ldots, 6\} 上の分布で、平均が m
に固定されたものの中でエントロピーを最大にする分布が p(k) ∝ e\^{}(θk)
の形になることをラグランジュの未定乗数法で示せ。また m = 3.5 のときの θ
を求めよ。

\hypertarget{ux554fux984cux2176-ux30afux30e9ux30e1ux30fcux30ebux30e9ux30aaux306eux4e0dux7b49ux5f0fux3068ux63a8ux5b9aux306eux52b9ux7387}{%
\subsection{問題Ⅶ　クラメール＝ラオの不等式と推定の効率}\label{ux554fux984cux2176-ux30afux30e9ux30e1ux30fcux30ebux30e9ux30aaux306eux4e0dux7b49ux5f0fux3068ux63a8ux5b9aux306eux52b9ux7387}}

次の文章を読んで、【　】に適した式・語句を、\{　\}
からは適切なものを一つ選びその番号を、それぞれ答えよ。また、問1・問2では指示にしたがって答えよ。

\begin{enumerate}
\def\labelenumi{(\arabic{enumi})}
\tightlist
\item
  データ x からパラメータ ξ を推定する関数 T(x) を推定量という。すべての
  ξ で E{[}T{]} = ξ
  を満たすものを不偏推定量という。E{[}T{\fallbackfont ⁱ}{]} = ∫
  T{\fallbackfont ⁱ}(x) p(x; ξ) dx = ξ{\fallbackfont ⁱ} の両辺を
  ξ{\fallbackfont ʲ} で微分し、問題Ⅰの E{[}∂{\fallbackfont ⱼ}l{]} = 0
  を使うと
\end{enumerate}

\[E\big[ (T^i - \xi^i)\, \partial_j l \big] = \text{【ア】}\]

を得る。推定誤差とスコアは「ちょうど単位行列ぶん」相関している。パラメータが1つのとき、コーシー＝シュワルツの不等式
(E{[}(T − ξ)∂l{]})² ≤ E{[}(T − ξ)²{]} E{[}(∂l)²{]} から

\[V[T] \ge \text{【イ】}\]

が従う（問1）。多次元では共分散行列について Cov(T) − g{\fallbackfont ⁻}¹
が半正定値になる。これがクラメール＝ラオの不等式であり、フィッシャー情報が大きい（分布の変化を見分けやすい）ほど、推定の精度の限界も良くなる。

独立同分布の N 個のデータでは、対数尤度が和になるのでフィッシャー情報は
【ウ】 となり、不偏推定量の分散は 1/N に比例して小さくできる限界をもつ。

\begin{enumerate}
\def\labelenumi{(\arabic{enumi})}
\setcounter{enumi}{1}
\tightlist
\item
  不等式の等号が成り立つ推定量を有効推定量という。コーシー＝シュワルツの等号条件から、有効であるためにはスコアが推定誤差に比例し、∂l
  = g(ξ)(T − ξ) の形でなければならない。これを ξ
  について積分すると、p(x; ξ) は T
  を統計量とする指数型分布族であり、しかも ξ はその 【エ】
  でなければならないことがわかる。つまり有効推定量が存在するのは、指数型分布族を期待値パラメータで表したときに限られ、そのとき有効推定量は統計量
  F の 【オ】 である。
\end{enumerate}

たとえばベルヌーイ分布の N
回の試行で成功割合を推定量とすると、その分散は p(1 − p)/N
である。問題Ⅰより g = 1/(p(1 − p)) なので、クラメール＝ラオの不等式で
【カ】 が成り立ち、成功割合は有効推定量である。

\begin{enumerate}
\def\labelenumi{(\arabic{enumi})}
\setcounter{enumi}{2}
\tightlist
\item
  一般のモデルでは有限の N で有効推定量は存在しないが、最尤推定量
  {\fallbackfont ξ̂} は N → ∞ で
\end{enumerate}

\[\sqrt{N}\, (\hat\xi - \xi) \longrightarrow \mathcal{N}\big(0,\ \text{【キ】}\big)\]

と分布収束し、漸近的に限界を達成する（漸近有効性）。さらに指数型分布族の中に曲がって埋め込まれたモデル（曲指数型分布族）では、推定誤差の
1/N² の項が、モデルが周囲の空間の中でどれだけ曲がっているかを表す 【ク】
で決まることが知られている。情報幾何学は、推定の精度を曲率という幾何学的な量で語ることを可能にした。

クラメール＝ラオの不等式の幾何学的な意味として、フィッシャー計量で測った距離が大きい2つの分布ほど
\{ ケ \}。

\begin{itemize}
\tightlist
\item
  ① データから識別しやすく、パラメータの推定精度も高い
\item
  ② データから識別しにくい
\item
  ③ 推定精度とは無関係である
\end{itemize}

問1　パラメータが1つの場合に、不等式 V{[}T{]} ≥ 【イ】
を導け。等号成立条件も述べよ。

問2　分散 σ² が既知の正規分布 N(μ, σ²) から N
個の独立なデータを得たとき、標本平均の分散がクラメール＝ラオの限界に一致することを確かめよ。

\hypertarget{ux554fux984cux2177-ux6b63ux898fux5206ux5e03ux65cfux306eux5e7eux4f55ux53ccux66f2ux5e73ux9762}{%
\subsection{問題Ⅷ　正規分布族の幾何（双曲平面）}\label{ux554fux984cux2177-ux6b63ux898fux5206ux5e03ux65cfux306eux5e7eux4f55ux53ccux66f2ux5e73ux9762}}

次の文章を読んで、【　】に適した式・数値・語句を、\{　\}
からは適切なものを一つ選びその番号を、それぞれ答えよ。また、問1・問2では指示にしたがって答えよ。

\begin{enumerate}
\def\labelenumi{(\arabic{enumi})}
\tightlist
\item
  1次元正規分布全体 \{N(μ, σ²) : μ ∈ {\fallbackfont ℝ}, σ \textgreater{}
  0\} を (μ, σ)
  を座標とする上半平面とみなす。問題Ⅰの結果から、フィッシャー計量は
\end{enumerate}

\[ds^2 = \frac{d\mu^2 + \text{【イ】}\, d\sigma^2}{\sigma^2}\]

である。μ = √2 u と変数変換すると ds² = 2(du² + dσ²)/σ² となる。(du² +
dσ²)/σ² はガウス曲率 −1
の双曲平面（ポアンカレの上半平面モデル）の計量であり、計量を定数 c
倍するとガウス曲率は 1/c 倍になるので、正規分布族のガウス曲率は K =
【ロ】
で一定である（問1）。正規分布の集まりは、負の定曲率をもつ双曲幾何の世界なのである。

\begin{enumerate}
\def\labelenumi{(\arabic{enumi})}
\setcounter{enumi}{1}
\tightlist
\item
  上半平面モデルの測地線は、境界 σ = 0
  に垂直な半直線と、中心が境界上にある半円である。(μ, σ)
  座標に戻すと、測地線は μ = 一定の半直線と、中心が μ 軸上にあり横方向に
  √2 倍引き伸ばされた 【ハ】 になる。
\end{enumerate}

μ = 一定の半直線に沿った距離は ∫ √2 dσ/σ で計算でき、N(μ,
σ{\fallbackfont ₁}²) と N(μ, σ{\fallbackfont ₂}²) の距離は 【ニ】
である。距離が σ の比だけで決まり、σ が小さい領域ほど同じ σ
の差でも遠く離れて見える。分散の小さい分布ほど、わずかな変化でも区別しやすいことの表れである。

では平均だけが異なる N(0, σ²) と N(Δ, σ²)
を結ぶ測地線はどうなるか。上の性質を踏まえると、この測地線は \{ ホ \}。

\begin{itemize}
\tightlist
\item
  ① 途中で σ が大きくなる方向にふくらんだ曲線になる
\item
  ② σ が一定の線分である
\item
  ③ 途中で σ が小さくなる方向にへこんだ曲線になる
\end{itemize}

\begin{enumerate}
\def\labelenumi{(\arabic{enumi})}
\setcounter{enumi}{2}
\tightlist
\item
  2つの正規分布の KL ダイバージェンスは
\end{enumerate}

\[D\big(\mathcal{N}(\mu_1, \sigma_1^2) \,\|\, \mathcal{N}(\mu_2, \sigma_2^2)\big) = \text{【ヘ】}\]

で与えられる（問2）。正規分布族は指数型分布族であり（問題Ⅲ）、θ 座標
(μ/σ², −1/(2σ²)) と η 座標 (μ, μ² + σ²) をもつ。したがって
e-接続・m-接続に関する曲率はともに 【ト】
である。同じ空間が、レヴィ＝チヴィタ接続では曲がっており（K =
【ロ】）、e-接続・m-接続ではまっすぐである。「平坦かどうか」は接続の選び方に依存するという、情報幾何学の要点がここに現れている。

\begin{enumerate}
\def\labelenumi{(\arabic{enumi})}
\setcounter{enumi}{3}
\tightlist
\item
  平均を固定した多変量正規分布 N(0, Σ)（Σ
  は正定値対称行列）のフィッシャー計量は
\end{enumerate}

\[ds^2 = \text{【チ】}\; \mathrm{tr}\big( \Sigma^{-1} d\Sigma\, \Sigma^{-1} d\Sigma \big)\]

で与えられる。1次元（Σ = σ²）では dΣ = 2σ dσ なので、(1) の σ 方向の成分
2dσ²/σ² と一致する。この計量は Σ {\fallbackfont ↦}
AΣA{\fallbackfont ᵀ}（A は正則行列）で不変であり、2点
Σ{\fallbackfont ₁}, Σ{\fallbackfont ₂} の測地距離は、行列 【リ】
の固有値 λ{\fallbackfont ₁}, \ldots, λ{\fallbackfont ₙ} を用いて d² =
(1/2) Σ{\fallbackfont ᵢ} (log λ{\fallbackfont ᵢ})²
と表される。共分散行列の比較、拡散テンソル画像の解析、脳波の信号処理などで実際に使われている距離である。

問1　問題Ⅳの微分幾何学の公式（F = 0 のとき）

\[K = -\frac{1}{2\sqrt{EG}} \left[ \frac{\partial}{\partial u}\left( \frac{G_u}{\sqrt{EG}} \right) + \frac{\partial}{\partial v}\left( \frac{E_v}{\sqrt{EG}} \right) \right]\]

を u = μ、v = σ、E = 1/σ²、G = 2/σ² に適用して、K = 【ロ】 を確かめよ。

問2　【ヘ】を導け。また μ{\fallbackfont ₂} = μ{\fallbackfont ₁} +
dμ、σ{\fallbackfont ₂} = σ{\fallbackfont ₁} + dσ
として2次まで展開し、問題Ⅱ(2) のとおり (1/2) ds²
が得られることを確かめよ。

\hypertarget{ux89e3ux7b54ux3068ux89e3ux8aacux554fux984cux2170ux2173}{%
\subsection{解答と解説（問題Ⅰ〜Ⅳ）}\label{ux89e3ux7b54ux3068ux89e3ux8aacux554fux984cux2170ux2173}}

空欄は表で、問は答案例で示す。答案例は要点を押さえた1つの書き方であり、これに限らない。

\hypertarget{ux554fux984cux2170}{%
\subsubsection{問題Ⅰ}\label{ux554fux984cux2170}}

\begin{longtable}[]{@{}lll@{}}
\toprule\noalign{}
空欄 & 解答 & ポイント \\
\midrule\noalign{}
\endhead
\bottomrule\noalign{}
\endlastfoot
ア & 0 & ∫p dx = 1 の微分 \\
イ & 0 & ∫∂{\fallbackfont ᵢ}∂{\fallbackfont ⱼ}p dx =
∂{\fallbackfont ᵢ}∂{\fallbackfont ⱼ}∫p dx \\
ウ & 1/(θ(1 − θ)) & V{[}x{]} = θ(1 − θ) \\
エ & 1/σ² & \\
オ & 2/σ² & z = (x − μ)/σ とすると E{[}(z² − 1)²{]} = 2 \\
カ & 4 & dθ/d{\fallbackfont ϕ} = 2 sin {\fallbackfont ϕ} cos
{\fallbackfont ϕ} \\
キ & π & ∫ 2 d{\fallbackfont ϕ}（0〜π/2） \\
ク & 2 & 4 = 2² \\
ケ & ① & \\
\end{longtable}

問1　∂{\fallbackfont ᵢ}∂{\fallbackfont ⱼ}l =
∂{\fallbackfont ᵢ}(∂{\fallbackfont ⱼ}p/p) =
(∂{\fallbackfont ᵢ}∂{\fallbackfont ⱼ}p)/p −
(∂{\fallbackfont ᵢ}p)(∂{\fallbackfont ⱼ}p)/p² =
(∂{\fallbackfont ᵢ}∂{\fallbackfont ⱼ}p)/p − ∂{\fallbackfont ᵢ}l
∂{\fallbackfont ⱼ}l。期待値をとると
E{[}(∂{\fallbackfont ᵢ}∂{\fallbackfont ⱼ}p)/p{]} =
∫∂{\fallbackfont ᵢ}∂{\fallbackfont ⱼ}p dx =
∂{\fallbackfont ᵢ}∂{\fallbackfont ⱼ}∫p dx = 0
なので、−E{[}∂{\fallbackfont ᵢ}∂{\fallbackfont ⱼ}l{]} =
E{[}∂{\fallbackfont ᵢ}l ∂{\fallbackfont ⱼ}l{]} = g{\fallbackfont ᵢⱼ}。

問2　l = −λ + x log λ − log x! より ∂l = x/λ − 1 = (x − λ)/λ。V{[}x{]} =
λ なので g(λ) = λ/λ² = 1/λ。λ = ν² なら dλ/dν = 2ν
で、{\fallbackfont g̃}(ν) = (2ν)² · (1/ν²) =
4（定数）。ベルヌーイと同じく「平方根をとると平坦になる」。

\hypertarget{ux554fux984cux2171}{%
\subsubsection{問題Ⅱ}\label{ux554fux984cux2171}}

\begin{longtable}[]{@{}lll@{}}
\toprule\noalign{}
空欄 & 解答 & ポイント \\
\midrule\noalign{}
\endhead
\bottomrule\noalign{}
\endlastfoot
イ & 0 & ∫q dx − ∫p dx \\
ロ & ① & \\
ハ & 0 & スコアの平均 \\
ニ & g{\fallbackfont ᵢⱼ} & 問題Ⅰ問1 \\
ホ & (μ{\fallbackfont ₁} − μ{\fallbackfont ₂})²/(2σ²) & 問2 \\
ヘ & λ{\fallbackfont ₁} log(λ{\fallbackfont ₁}/λ{\fallbackfont ₂}) −
λ{\fallbackfont ₁} + λ{\fallbackfont ₂} & E{[}x{]} =
λ{\fallbackfont ₁} \\
ト & 1/λ{\fallbackfont ₁} & 問題Ⅰ問2の g と一致 \\
チ & −log t & p f(q/p) = p log(p/q) \\
リ & 1 & f″(t) = 1/t² \\
ヌ & 1/2 & f″(t) = (1/2) t\^{}(−3/2) \\
\end{longtable}

問1　log t ≤ t − 1（等号は t = 1 のみ）に t = q(x)/p(x) を代入して p(x)
を掛け、p \textgreater{} 0 の範囲で積分すると −D(p‖q) ≤ ∫q dx − 1 ≤
0。等号には、p \textgreater{} 0 のほとんどすべての x で q/p = 1
であることが必要十分で、そのとき p = q。

問2　log(p{\fallbackfont ₁}/p{\fallbackfont ₂}) = {[}(x −
μ{\fallbackfont ₂})² − (x −
μ{\fallbackfont ₁})²{]}/(2σ²)。p{\fallbackfont ₁} のもとで E{[}(x −
μ{\fallbackfont ₁})²{]} = σ²、E{[}(x − μ{\fallbackfont ₂})²{]} = σ² +
(μ{\fallbackfont ₁} − μ{\fallbackfont ₂})² なので、D =
(μ{\fallbackfont ₁} − μ{\fallbackfont ₂})²/(2σ²)。これは
(1/2)·(1/σ²)·(μ{\fallbackfont ₁} − μ{\fallbackfont ₂})²
と書け、局所展開がここでは正確に成り立つ。

\hypertarget{ux554fux984cux2172}{%
\subsubsection{問題Ⅲ}\label{ux554fux984cux2172}}

\begin{longtable}[]{@{}
  >{\raggedright\arraybackslash}p{(\columnwidth - 4\tabcolsep) * \real{0.3333}}
  >{\raggedright\arraybackslash}p{(\columnwidth - 4\tabcolsep) * \real{0.3333}}
  >{\raggedright\arraybackslash}p{(\columnwidth - 4\tabcolsep) * \real{0.3333}}@{}}
\toprule\noalign{}
\begin{minipage}[b]{\linewidth}\raggedright
空欄
\end{minipage} & \begin{minipage}[b]{\linewidth}\raggedright
解答
\end{minipage} & \begin{minipage}[b]{\linewidth}\raggedright
ポイント
\end{minipage} \\
\midrule\noalign{}
\endhead
\bottomrule\noalign{}
\endlastfoot
あ & g{\fallbackfont ᵢⱼ}（フィッシャー情報） & \\
い & log(1 + e\^{}θ) & Σ（x = 0, 1）e\^{}(θx) \\
う & e\^{}θ/(1 + e\^{}θ) & ロジットの逆関数 \\
え & p(1 − p) & \\
お & −1/(2σ²) & x² の係数 \\
か & μ² + σ² & \\
き & 負のエントロピー（エントロピーの符号を変えたもの） &
基準測度に関するもの \\
く & 0 & フェンシェルの等式 \\
け & η log η + (1 − η) log(1 − η) & \\
こ & ① & 尤度方程式 ∂{\fallbackfont ᵢ}ψ({\fallbackfont θ̂}) = 標本平均 \\
\end{longtable}

問1　e\^{}ψ = ∫exp(θ · F) dx を θ{\fallbackfont ⁱ} で微分すると
∂{\fallbackfont ᵢ}ψ e\^{}ψ = ∫F{\fallbackfont ᵢ} exp(θ · F) dx、両辺を
e\^{}ψ で割って ∂{\fallbackfont ᵢ}ψ = E{[}F{\fallbackfont ᵢ}{]}。さらに
θ{\fallbackfont ʲ} で微分すると ∂{\fallbackfont ⱼ}∂{\fallbackfont ᵢ}ψ =
∫F{\fallbackfont ᵢ} (F{\fallbackfont ⱼ} − ∂{\fallbackfont ⱼ}ψ) p dx =
E{[}F{\fallbackfont ᵢ}F{\fallbackfont ⱼ}{]} −
E{[}F{\fallbackfont ᵢ}{]}E{[}F{\fallbackfont ⱼ}{]} =
Cov(F{\fallbackfont ᵢ}, F{\fallbackfont ⱼ})。

問2　p(x; λ) = (1/x!) exp(x log λ − λ) なので、基準測度を 1/x! として θ
= log λ、F(x) = x、ψ(θ) = e\^{}θ。η = ψ′(θ) = e\^{}θ = λ、θ
座標でのフィッシャー情報 ψ″(θ) = e\^{}θ = λ（λ 座標の 1/λ に (dλ/dθ)² =
λ² を掛けたものと一致）。φ(η) = θη − ψ = η log η − η。

\hypertarget{ux554fux984cux2173}{%
\subsubsection{問題Ⅳ}\label{ux554fux984cux2173}}

\begin{longtable}[]{@{}lll@{}}
\toprule\noalign{}
空欄 & 解答 & ポイント \\
\midrule\noalign{}
\endhead
\bottomrule\noalign{}
\endlastfoot
ア & α/2 & Γ\^{}(0) = E{[}∂{\fallbackfont ᵢ}∂{\fallbackfont ⱼ}l
∂{\fallbackfont ₖ}l{]} + T/2 \\
イ & −α & \\
ウ & 0 & \\
エ & 0 & \\
オ & 0 & \\
カ & 平坦（双対平坦） & \\
キ & (1 − t)p + tq & \\
ク & log ∫ p\^{}(1−t) q\^{}t dx & \\
ケ & p(1 − p)(1 − 2p) & ψ″ = p(1 − p) を θ で微分、dp/dθ = p(1 − p) \\
コ & ① & \\
\end{longtable}

問1　g{\fallbackfont ᵢⱼ} = ∫∂{\fallbackfont ᵢ}l ∂{\fallbackfont ⱼ}l p dx
を ξ{\fallbackfont ᵏ} で微分すると、∂{\fallbackfont ₖ}p = p
∂{\fallbackfont ₖ}l より

\[\partial_k g_{ij} = E[\partial_k\partial_i l\, \partial_j l] + E[\partial_i l\, \partial_k\partial_j l] + E[\partial_i l\, \partial_j l\, \partial_k l]\]

一方、Γ\^{}(e){\fallbackfont ₖᵢ},{\fallbackfont ⱼ} =
E{[}∂{\fallbackfont ₖ}∂{\fallbackfont ᵢ}l
∂{\fallbackfont ⱼ}l{]}、Γ\^{}(m){\fallbackfont ₖⱼ},{\fallbackfont ᵢ} =
E{[}(∂{\fallbackfont ₖ}∂{\fallbackfont ⱼ}l + ∂{\fallbackfont ₖ}l
∂{\fallbackfont ⱼ}l) ∂{\fallbackfont ᵢ}l{]}
であり、両者の和は右辺の3項にちょうど一致する。

問2　log p(x; μ) = −(x − μ)²/2 + 定数 なので、(1 − t) log
p{\fallbackfont ₁} + t log p{\fallbackfont ₂} = −{[}(1 − t)(x −
μ{\fallbackfont ₁})² + t(x − μ{\fallbackfont ₂})²{]}/2 + 定数 = −(x −
m{\fallbackfont ₜ})²/2 + (x によらない項)、ただし m{\fallbackfont ₜ} =
(1 − t)μ{\fallbackfont ₁} + tμ{\fallbackfont ₂}。正規化すると
N(m{\fallbackfont ₜ}, 1)。指数型分布族は
e-平坦なので、e-測地線は族の中にとどまり、θ 座標（ここでは
μ）で直線になる。

\hypertarget{ux89e3ux7b54ux3068ux89e3ux8aacux554fux984cux2174ux2177}{%
\subsection{解答と解説（問題Ⅴ〜Ⅷ）}\label{ux89e3ux7b54ux3068ux89e3ux8aacux554fux984cux2174ux2177}}

\hypertarget{ux554fux984cux2174}{%
\subsubsection{問題Ⅴ}\label{ux554fux984cux2174}}

\begin{longtable}[]{@{}lll@{}}
\toprule\noalign{}
空欄 & 解答 & ポイント \\
\midrule\noalign{}
\endhead
\bottomrule\noalign{}
\endlastfoot
イ & D(q‖p)（Q の分布から P の分布への KL） & 問題Ⅲの ψ, φ を代入 \\
ロ & ∇ψ(θ\_Q) · (θ\_P − θ\_Q) & η\_Q = ∇ψ(θ\_Q) \\
ハ & 凸性 & 凸関数は接平面より上にある \\
ニ & \textbar θ\_P − θ\_Q\textbar²/2 & \\
ホ & (θ\_P − θ\_Q) · (η\_R − η\_Q) & 問1 \\
ヘ & Σ{\fallbackfont ᵢ} X{\fallbackfont ⁱ}Y{\fallbackfont ᵢ} &
∂η{\fallbackfont ᵢ}/∂θ{\fallbackfont ʲ} = g{\fallbackfont ᵢⱼ} による \\
ト & D(r‖q) & \\
チ & ① & \\
\end{longtable}

イの確認：p, q を P, Q の分布とすると、D(q‖p) = E\_q{[}(θ\_Q − θ\_P) · F
− ψ(θ\_Q) + ψ(θ\_P){]} = (θ\_Q − θ\_P) · η\_Q − ψ(θ\_Q) + ψ(θ\_P) =
ψ(θ\_P) + φ(η\_Q) − θ\_P · η\_Q。

問1　D{[}P:Q{]} + D{[}Q:R{]} − D{[}P:R{]} に定義を代入すると、ψ(θ\_P) と
φ(η\_R) が消え、ψ(θ\_Q) + φ(η\_Q) = θ\_Q · η\_Q を用いて

\[D[P:Q] + D[Q:R] - D[P:R] = \theta_Q\cdot\eta_Q - \theta_P\cdot\eta_Q - \theta_Q\cdot\eta_R + \theta_P\cdot\eta_R = (\theta_P - \theta_Q)\cdot(\eta_R - \eta_Q)\]

問2　ψ が狭義凸なら ψ(θ\_P) ≥ ψ(θ\_Q) + ∇ψ(θ\_Q) · (θ\_P − θ\_Q)
で、等号は θ\_P = θ\_Q のときに限る。左辺から右辺を引いたものが
D{[}P:Q{]} である。

\hypertarget{ux554fux984cux2175}{%
\subsubsection{問題Ⅵ}\label{ux554fux984cux2175}}

\begin{longtable}[]{@{}lll@{}}
\toprule\noalign{}
空欄 & 解答 & ポイント \\
\midrule\noalign{}
\endhead
\bottomrule\noalign{}
\endlastfoot
あ & (1/N) Σ{\fallbackfont ₖ} log p(x{\fallbackfont ₖ};
θ)（平均対数尤度） & \\
い & m-射影 & \\
う & e-平坦 & \\
え & k/N & 成功割合 \\
お & exp(θ · F(x) − ψ(θ)) & \\
か & ラグランジュの未定乗数 & \\
き & e-射影 & 第1引数を動かす \\
く & 増加しない（単調非増加） & 各ステップが最小化 \\
け & 局所的な最小点（停留点） & 初期値に依存する \\
こ & ① & \\
\end{longtable}

問1　D({\fallbackfont p̂}‖p\_θ) = Σ {\fallbackfont p̂} log
{\fallbackfont p̂} − Σ{\fallbackfont ₓ} {\fallbackfont p̂}(x) log p(x;
θ)。{\fallbackfont p̂}(x) は値 x の出現割合なので、第2項は (1/N)
Σ{\fallbackfont ₖ} log p(x{\fallbackfont ₖ}; θ) に等しい。第1項は θ
によらないから、D を最小にする θ は平均対数尤度を最大にする
θ（最尤推定量）である。これは {\fallbackfont p̂} から M 上で
D({\fallbackfont p̂}‖q) を最小にする点を探すこと、すなわち m-射影である。

問2　L = −Σ p{\fallbackfont ₖ} log p{\fallbackfont ₖ} −
λ{\fallbackfont ₀}(Σ p{\fallbackfont ₖ} − 1) − λ{\fallbackfont ₁}(Σ k
p{\fallbackfont ₖ} − m) を p{\fallbackfont ₖ} で微分して 0 とおくと −log
p{\fallbackfont ₖ} − 1 − λ{\fallbackfont ₀} − λ{\fallbackfont ₁}k =
0、よって p{\fallbackfont ₖ} ∝ e\^{}(θk)（θ =
−λ{\fallbackfont ₁}）。平均は θ について単調増加なので、m ごとに θ
は一意。一様分布は平均 3.5 を満たしエントロピーが最大なので、m = 3.5
のとき θ = 0。

\hypertarget{ux554fux984cux2176}{%
\subsubsection{問題Ⅶ}\label{ux554fux984cux2176}}

\begin{longtable}[]{@{}lll@{}}
\toprule\noalign{}
空欄 & 解答 & ポイント \\
\midrule\noalign{}
\endhead
\bottomrule\noalign{}
\endlastfoot
ア & δ{\fallbackfont ⁱⱼ}（i = j で 1、それ以外 0） &
∂{\fallbackfont ⱼ}ξ{\fallbackfont ⁱ} \\
イ & 1/g(ξ) & \\
ウ & N g(ξ) & 独立なら情報は足し算 \\
エ & 期待値パラメータ（η 座標） & \\
オ & 標本平均 (1/N) Σ F(x{\fallbackfont ₖ}) & \\
カ & 等号 & p(1 − p)/N = 1/(N g) \\
キ & g(ξ){\fallbackfont ⁻}¹ & \\
ク & 曲率（埋め込み曲率、e-曲率） & エフロンの統計的曲率 \\
ケ & ① & \\
\end{longtable}

問1　アで i = j とすると E{[}(T − ξ) ∂l{]} =
1。コーシー＝シュワルツの不等式より 1 ≤ E{[}(T − ξ)²{]} E{[}(∂l)²{]} =
V{[}T{]} g(ξ)。等号は ∂l = c(ξ)(T − ξ)
がほとんど確実に成り立つときで、両辺に (T − ξ) を掛けて期待値をとると
c(ξ) = g(ξ)。

問2　1個あたり g = 1/σ²、N 個で N/σ² なので限界は σ²/N。標本平均の分散は
σ²/N で一致する。正規分布は x を統計量、μ
を期待値パラメータとする指数型分布族なので、(2) の結論とも整合する。

\hypertarget{ux554fux984cux2177}{%
\subsubsection{問題Ⅷ}\label{ux554fux984cux2177}}

\begin{longtable}[]{@{}lll@{}}
\toprule\noalign{}
空欄 & 解答 & ポイント \\
\midrule\noalign{}
\endhead
\bottomrule\noalign{}
\endlastfoot
イ & 2 & 問題Ⅰ(2) \\
ロ & −1/2 & 問1 \\
ハ & 半楕円 & 半円を μ 方向に √2 倍 \\
ニ & √2 \textbar log(σ{\fallbackfont ₂}/σ{\fallbackfont ₁})\textbar{}
& \\
ホ & ① & \\
ヘ & log(σ{\fallbackfont ₂}/σ{\fallbackfont ₁}) + (σ{\fallbackfont ₁}² +
(μ{\fallbackfont ₁} − μ{\fallbackfont ₂})²)/(2σ{\fallbackfont ₂}²) − 1/2
& 問2 \\
ト & 0 & \\
チ & 1/2 & \\
リ & Σ{\fallbackfont ₁⁻}¹Σ{\fallbackfont ₂} & \\
\end{longtable}

ホの理由：μ 方向に dμ 動くコストは dμ/σ に比例するので、σ
の大きい領域を通って横に移動するほうが短くなる。σ
一定の線分は測地線ではなく、測地線は上に（σ
の大きい方へ）ふくらむ半楕円である。

問1　√(EG) = √2/σ²。G は u によらないので第1項は 0。E{\fallbackfont ᵥ} =
−2/σ³ より E{\fallbackfont ᵥ}/√(EG) = −√2/σ、これを v（= σ）で微分すると
√2/σ²。よって K = −(σ²/(2√2)) · (√2/σ²) = −1/2。

問2　log(p{\fallbackfont ₁}/p{\fallbackfont ₂}) =
log(σ{\fallbackfont ₂}/σ{\fallbackfont ₁}) − (x −
μ{\fallbackfont ₁})²/(2σ{\fallbackfont ₁}²) + (x −
μ{\fallbackfont ₂})²/(2σ{\fallbackfont ₂}²)。p{\fallbackfont ₁} のもとで
E{[}(x − μ{\fallbackfont ₁})²{]} = σ{\fallbackfont ₁}²、E{[}(x −
μ{\fallbackfont ₂})²{]} = σ{\fallbackfont ₁}² + (μ{\fallbackfont ₁} −
μ{\fallbackfont ₂})² なので【ヘ】を得る。σ{\fallbackfont ₂} = σ +
dσ、μ{\fallbackfont ₂} = μ + dμ（σ{\fallbackfont ₁} =
σ、μ{\fallbackfont ₁} = μ）とおいて2次まで展開すると、log(1 + dσ/σ) ≈
dσ/σ − dσ²/(2σ²)、(σ² + dμ²)/(2(σ + dσ)²) ≈ 1/2 − dσ/σ + 3dσ²/(2σ²) +
dμ²/(2σ²) より、D ≈ (dμ² + 2dσ²)/(2σ²) = (1/2) ds²。

\end{document}
